% !PRINTED 380
\weberpage{380}
%<<CONT

Eine complexe Zahl $\omega = \xi + i\eta$ heisst bei positivem $\eta$ \gesperrt{reducirt}, wenn
\weq{2}{-\tfrac12 \leqq \xi \leqq \tfrac12, \quad \xi^2+\eta^2 \geqq 1.}
Betrachtet man $\xi$, $\eta$ als rechtwinklige Coordinaten in einer Ebene, so wird jede Zahl $\omega$ durch einen Punkt dieser Ebene veranschaulicht, und die Lage der Punkte, die den reducirten Zahlen entsprechen, wird durch das in der Fig. 26 schraffirte Feld (mit Einschluss der Grenzen), das wir das Grundfeld nennen, veranschaulicht.

% SOURCE FIGURE CROP: Bd. I registered scan 406 (printed p. 380), 600 dpi,
% crop x=350,y=1005,w=1180,h=900; source-faithful raster, Fig. 26.
\weberfigureimage{weber-b1-fig26.png}{0.48}

Aus den Bedingungen (2) folgt dann noch, was auch in der Figur leicht zu bestätigen ist,
\weq{3}{\eta^2 \geqq \tfrac34.}
Geht man nun von einer Zahl $\omega$ zu einer äquivalenten Zahl $\omega_1$ über, so ergiebt die Formel (16) des vorigen Paragraphen
\weq{4}{\omega_1 = \frac{(\alpha\xi+\beta)+i\alpha\eta}{(\gamma\xi+\delta)+i\gamma\eta} = \frac{[(\alpha\xi+\beta)(\gamma\xi+\delta)+\alpha\gamma\eta^2]+i\varepsilon\eta}{(\gamma\xi+\delta)^2+\gamma^2\eta^2},}
und diese Formel zeigt, dass $\eta_1$ dasselbe oder das entgegengesetzte Zeichen wie $\eta$ hat, je nachdem $\varepsilon = +1$ oder $=-1$ ist, also je nachdem die Aequivalenz eine eigentliche oder eine uneigentliche ist.

Beschränken wir uns also auf Zahlen mit positiv imaginärem Bestandtheil, so kommt nur die eigentliche Aequivalenz in Betracht, und wir beweisen jetzt den Fundamentalsatz:

1. \gesperrt{Jede imaginäre quadratische Irrationalzahl $x+y\sqrt{d}$ mit positiv imaginärem Bestandtheil ist mit einer reducirten Zahl äquivalent.}

Verstehen wir unter $\alpha$ die dem $\xi$ nächstgelegene ganze Zahl, so wird $\omega-\alpha$ der ersten der Bedingungen (2) genügen, dass nämlich $-\tfrac12 \leqq \xi-\alpha \leqq \tfrac12$ ist.

Wenn nun
\weq{5}{r^2 = (\xi-\alpha)^2+\eta^2 = (\omega-\alpha)(\omega'-\alpha)}
grösser oder gleich 1 ist, so ist $\omega-\alpha$, was mit $\omega$ äquivalent ist, bereits reducirt; anderenfalls setzen wir, indem wir eine Kettenbruchentwickelung anwenden,
%CONT>>
